\section{The calculus, the changed commands, and the PA equivalence}
\label{cal:section}

This section gives the ordinary mathematical proof for the literal
machine $M_{194C}$.  The numerical calculus and its diagonal calibration
come from the supplied PA194 research packet; the changed command protocol
comes from the supplied PA195 Combined construction.  The new result is
their composition with the 194-state table and the fresh source and tape
certificates.  The detailed inherited arguments are retained in
\file{construction/proof_sources/}; in particular,
\file{original_pa195_calibration.tex},
\file{diagonal_calculus.tex}, and \file{combined_commands.tex} identify the
proof interfaces being reused.  The present argument is a proof of a
PA-provability statement.  An end-to-end Lean or Rocq theorem for
$M_{194C}$ is not claimed.

\subsection{A completely specified numerical calculus}
\label{cal:calculus}

Let
\[
 p(a,b)=\frac{(a+b)(a+b+1)}2+b,
 \qquad
 \begin{array}{ll}
 I(A,B)=4p(A,B),&M(a,b)=4p(a,b)+1,\\
 U(a,B)=4p(B,a)+2,&E(a,b)=4p(a,b)+3.
 \end{array}
\]
Every natural number is a formula code.  Code $0$ is truth; code
$1=M(0,0)$ is the designated contradiction.  The intended readings of
$I,M,U,E$ are implication, membership, universal quantification, and
equality.  Notice the transposition in $U$.  Write
\[
 N(A)=I(A,1),\qquad
 A\wedge B=N(I(A,N(B))),\qquad
 \exists_a B=N(U(a,N(B))).
\]
In displays of coded formulas, $A\to B$ abbreviates $I(A,B)$ and
$\neg A$ abbreviates $N(A)$.  These are operations on natural numbers,
not additional primitive syntax.  The exceptional identity
$I(0,0)=0$ will be respected throughout.

For arbitrary natural parameters $D,A,H$, the six head schemas are
\begin{align*}
 L_1(D,A,H)&=(\exists_D U(D,H))\to H,\\
 L_2(D,A,H)&=U(D,U(A,H))\to U(A,U(D,H)),\\
 L_3(D,A,H)&=M(H,D)\to(E(H,A)\to M(A,D)),\\
 L_4(D,A,H)&=E(H,D)\to(M(A,H)\to M(A,D)),\\
 L_5(D,A,H)&=\exists_D E(H,D),\\
 L_6(D,A,H)&=U(D,A)\to\bigl(U(D,A\to H)\to U(D,H)\bigr).
\end{align*}
There is also the restricted implicational schema
\[
 \rho(a,b,c,u)=((c\to u)\to a)\to((a\to c)\to(b\to c)),
 \qquad u\le a,
\]
and the unconditional diagonal schema
\[
 \Gamma(B)=B\to U(B,B).
\]
Define the pointed-separation expressions
\begin{align*}
 V(a,b,c)&=M(a,b)\to U(c,M(a,b)),\\
 P(a,b,c)&=c\to M(c,b),\\
 Z(a,b,c)&=U\bigl(b,
     (M(a,b)\to P(a,b,c))\wedge(M(c,b)\to c)\bigr),\\
 Q(a,b,c)&=U\bigl(c,Z(a,b,c)\to\neg P(a,b,c)\bigr),\\
 G(a,b,c)&=\neg\bigl(V(a,b,c)\to Q(a,b,c)\bigr).
\end{align*}
The gate schema is $G(a,b,c)$ for $c\ne a,b$.  There are exactly two
inference rules:
\[
 \frac{A\qquad I(A,B)}{B}\quad\text{and}\quad
 \frac{B}{U(a,B)}\quad(a\in\mathbb N).
\]
The second is generalization of a theorem, not generalization under
arbitrary undischarged assumptions.  These schemas and rules define
$\mathcal C_\Gamma$.  No equality symmetry, equality substitution,
arbitrary formula substitution, or propositional-tautology oracle is
primitive.  Truth has a finite derivation because
$\rho(0,0,0,0)=0$ satisfies the dominance condition.

The retained PA196 implicational development supplies uniform finite
derivations of the classical implicational templates from the restricted
$\rho$ schema.  Thus the propositional reasoning below abbreviates
explicit proof compositions; it does not add a rule to
$\mathcal C_\Gamma$.  The relevant inherited source is
\file{PA196/Logic/ImplicationBasis.lean} in the pinned upstream snapshot,
with its arithmetic-certificate counterpart.

\subsection{Why this calculus has the consistency strength of PA}
\label{cal:logical}

We state both syntactic directions, since external truth in a structure
would not establish the requested internal equivalence.  Fix a standard
recursive proof predicate $\operatorname{Prf}_{\PA}$ for first-order PA
with full induction, and a recursive predicate
$\operatorname{Prf}_\Gamma(d,B)$ for the finite derivation trees just
specified.

\begin{lemma}[Uniform logical calibration]
\label{cal:logical-lemma}
There are primitive-recursive proof transformations $F,J$ whose
correctness is provable in PA and which satisfy
\begin{align*}
 \PA&\vdash\forall p\,
  \bigl(\operatorname{Prf}_{\PA}(p,\ulcorner0=1\urcorner)
       \to\operatorname{Prf}_\Gamma(F(p),1)\bigr),\\
 \PA&\vdash\forall d\,
  \bigl(\operatorname{Prf}_\Gamma(d,1)
       \to\operatorname{Prf}_{\PA}(J(d),\ulcorner0=1\urcorner)\bigr).
\end{align*}
\end{lemma}

\begin{proof}
For the forward direction, the inherited PA196 refutation compiler uses
primitive equality through three derived interfaces: two-way membership
rewriting, neighborhood equality, and an equality witness.  The only
uses of its removed equality-vacuity schema in the relevant logical and
arithmetic-certificate derivations are the reverse-rewriting and
neighborhood modules.  Here are the replacements, including the
freshness condition that makes them legitimate.

Every variable label occurring in a decoded formula of code $B$ is
strictly below $B$.  Indeed, the fields of any nonzero tagged constructor
are strictly smaller than its code; induction through its formula
children proves the claim, and code zero has no labels.  In particular
$x_B$ does not occur in its own decoded formula.  The gate supplies both
membership vacuity $V(a,b,c)$ and $\neg Q(a,b,c)$ by the classical
implicational templates.

For reverse rewriting, first let $a\ne b$ and take
$\chi(t)=M(d,t)$ or $M(t,d)$ with $d\ne a$, or take
$\chi(t)=M(t,t)$.  Put
\[
 C=E(a,b)\to(\chi(b)\to\chi(a)),\qquad
 \varphi=\neg C,\qquad c=\ulcorner\varphi\urcorner.
\]
The code bounds make $c$ distinct from every displayed variable label.
Apply the gate at $(a,a,c)$ and abbreviate its $Z,P,Q$ accordingly.
Its second separation conjunct and the propositional consequence
$\varphi\to\neg\chi(a)$ give
\[
 Z\to U(a,M(c,a)\to\neg\chi(a)).
\]
The equality witness, direct membership rewriting, and membership-only
vacuity yield the negative-specialization consequence
\[
 Z\to(M(c,b)\to\neg\chi(b)).
\]
More explicitly, this is the retained specialization using the witness
$\exists_a E(b,a)$ and the vacuity of
$M(c,b)\to\neg\chi(b)$ at $a$.  Its side conditions are exactly
$a\ne b,c$ and, in the nondiagonal cases, $a\ne d$; it needs no
equality-vacuity premise.  Under $\varphi$, $P$ supplies $M(c,a)$ and
$E(a,b)$ supplies $M(c,b)$, contradicting the other consequence
$\chi(b)$ of $\varphi$.  Hence
$\varphi\to(Z\to\neg P)$.  Generalization at $c$, quantifier
distribution, and $\Gamma(\varphi)$ give $\varphi\to Q$.
The gate's $\neg Q$ and double-negation elimination prove $C$.
For $a=b$ the result is propositional.  If the fixed label $d=a$,
direct rewriting first reaches $M(b,b)$, and the already proved diagonal
case reaches $M(a,a)$.  Thus both membership arguments can be rewritten
in either direction, without any remaining label restriction.

For neighborhood equality, let
\[
 R(a,b,z)=U(z,M(a,z)\leftrightarrow M(b,z))
          \wedge U(z,M(z,a)\leftrightarrow M(z,b)),
\]
where the biconditionals are the corresponding coded conjunctions.
At $z_0=\ulcorner E(a,b)\urcorner$, the diagonal axiom
$E(a,b)\to U(z_0,E(a,b))$ and the rewrites give
$E(a,b)\to R(a,b,z_0)$.  For $z\ne a,b,z_0$, let $B(t)$ be either
membership-only matrix in $R$.  The rewrites, the witness
$\exists_{z_0}E(z,z_0)$, and vacuity of $B(z)$ at $z_0$ give
\[
 U(z_0,B(z_0))\to B(z).
\]
Its antecedent is vacuous at $z$: membership vacuity extends through
the Boolean connectives and through the outer quantifier by $L_2$.
Generalization and distribution therefore give
$U(z_0,B(z_0))\to U(z,B(z))$.  The case $z=z_0$ is already done.
This proves $E(a,b)\to R(a,b,z)$ whenever $z\ne a,b$; existential
monotonicity and $L_5$ give the required neighborhood witness.  These
are precisely the three interfaces of the inherited compiler.
Replacing its two affected proof templates gives $F$.
All code choices, label comparisons, and proof concatenations above
are primitive recursive; the inherited arithmetic proof construction
checks the replacements by the same finite syntactic inductions.

For the reverse direction, work in a fixed conservative definitional
extension of PA with $2^z$, and set
\[
 \operatorname{Bit}(s,z)\equiv
 \exists u,v\,[s=(2u+1)2^z+v\ \wedge\ v<2^z].
\]
Translate $M(a,b)$ to $\operatorname{Bit}(x_a,x_b)$, $E(a,b)$ to
$x_a=x_b$, $U(a,B)$ to $\forall x_a\,T(B)$, and $0$ to $0=0$.
For implication, the constructor property is the provable equivalence
\[
 T(I(A,B))\leftrightarrow(T(A)\to T(B));
\]
when $A=B=0$, it uses the fixed truth equivalence rather than a false
assertion of syntactic constructor disjointness.  Bound auxiliary
variables in $\operatorname{Bit}$ are renamed before substitution.
PA proves $\operatorname{Bit}(s,z)\to z<s$, hence
$\neg T(1)=\neg\operatorname{Bit}(x_0,x_0)$.

The six logical heads have uniform first-order proof templates, $\rho$
is a classical tautology, and the code-freshness argument supplies
$T(\Gamma(B))$.  The substantive axiom supplier is the gate.  For any
arithmetic formula $\theta(z,\vec v)$, PA proves
\begin{equation}
 \forall n\,\exists s<2^n\,\forall z<n\,
  \bigl(\operatorname{Bit}(s,z)\leftrightarrow\theta(z,\vec v)\bigr).
 \label{cal:comprehension}
\end{equation}
Induct on $n$: use $s=0$ at zero, and retain $s$ or replace it by
$s+2^n$ according as $\theta(n,\vec v)$ fails or holds.  This is one
fixed proof pattern with a full PA induction instance.  Inserting an
arbitrary formula code into that pattern and checking the resulting
proof are primitive-recursive syntax operations; PA verifies their
correctness.  This supplies a uniform proof generator, not a uniform
truth predicate for arithmetic.

For the gate at $c\ne a,b$, put $\phi=T(c)$.  If $a\ne b$, use
\eqref{cal:comprehension} with bound $x_a+x_b+1$ and matrix
\[
  (\operatorname{Bit}(x_a,z)\lor z=x_b)\wedge\phi[z/x_b].
\]
Its witness gives $T(Z)\wedge T(P)$, and thus $\neg T(Q)$.
If the labels $a=b$ coincide, the first antecedent inside the binder
is $\operatorname{Bit}(z,z)$ and is false.  The correct matrix is then
$(z=x_b)\wedge\phi[z/x_b]$, with bound $x_b+1$.
This separate case prevents treating a bound occurrence of $x_a$ as
free.  Code freshness ensures that the separator variable $x_c$ does
not occur in $\phi$.  The membership-vacuity clause has its ordinary
first-order proof, so $T(G)$ follows.

Recursively translate a numerical derivation: insert these axiom
suppliers, concatenate at modus ponens, and append generalization at
its corresponding rule.  PA verifies the proof-code invariant by
induction on the derivation.  At conclusion $1$, append the fixed proof
of $\neg T(1)$ and eliminate the fixed definitional extension.  This
gives $J$.  The induction is on syntax and proof checking, so its
quantified assertion applies to internal, possibly nonstandard proof
codes as required.
\end{proof}

\subsection{Exact semantics of the Combined commands}
\label{cal:commands}

Write $[a;w]=p(a,w)$ and $[a_1,\ldots,a_k;w]$ for repeated pairing
with $a_1$ first.  The total Cantor projections obey
\[
 \operatorname{hd}[a;w]=a,\quad
 \operatorname{tl}[a;w]=w,\quad
 \operatorname{hd}(0)=\operatorname{tl}(0)=0.
\]
In particular $[0;0]=0$.  All following identities are identities of
natural numbers; they do not assume a free list constructor.

A command store is $\langle t,w,d,h\rangle$, with two visible formula
codes $t,h$, a stack $w$, and an offset $d$.  For $y=d+x$, the new
commands are
\begin{align*}
 \mathsf S\langle t,w,d,h\rangle
 &=\begin{cases}
     \mathsf{halt},&h=1,\\
     \langle h,\operatorname{tl}(w),d,\operatorname{hd}(w)\rangle,&h\ne1,
   \end{cases}\\
 \mathsf{MP}_x\langle t,w,d,h\rangle
 &=\begin{cases}
     \langle0,w,0,y\rangle,&I(t,y)=h,\\
     \langle0,w,y,0\rangle,&I(t,y)\ne h.
   \end{cases}
\end{align*}
Only $\mathsf S$ can halt.  Neither branch of MP pushes an entry.

Define $K(a,b,c)=G(a,b,c)$ when $c\ne a,b$ and $K(a,b,c)=0$
otherwise.  For $D=d+x$, $H=h+y$, and $A=c$, AX is exactly
\[
 \mathsf{AX}_{c,x,y}\langle t,w,d,h\rangle
 =\langle K(H,D,c),[q_1,\ldots,q_9,t;w],0,0\rangle,
\]
where, writing $L_i=L_i(D,A,H)$,
\[
 (q_1,\ldots,q_9)=
 \bigl(\Gamma(A),L_6,L_5,L_4,L_3,L_2,L_1,
              \rho(D,A,H,d),U(d,h)\bigr).
\]
The total number of pushed entries is ten, including the saved top.
The first eight generated entries and the gate have residue zero modulo
four; the last entry has residue two.  None is code one.  The separately
pushed $\Gamma(A)$ is present even when the gate guard fails.

\subsection{A proof compiler valid for every incoming stack}
\label{cal:compiler}

Command words execute from left to right.  The following exact
identities hold for every incoming store, including $t=1$:
\begin{align*}
 \mathsf{AX}_{c,x,y};\mathsf S^9:
 \langle t,w,d,h\rangle
 &\longmapsto\langle\rho(d+x,c,h+y,d),[t;w],0,U(d,h)\rangle,\\
 \mathsf{AX}_{c,x,y};\mathsf S^{10}:
 \langle t,w,d,h\rangle
 &\longmapsto\langle U(d,h),w,0,t\rangle.
\end{align*}
The first Shift tests the cache cleared by AX; the next nine test the
generated entries, never the saved top.  The residue calculation proves
that all these Shifts continue.

Define
\begin{align*}
 \mathsf{Push}&=\mathsf{AX}_{0,0,0};\mathsf S^9;\mathsf{MP}_0,\\
 \mathsf{Pop}&=\mathsf{AX}_{0,0,0};\mathsf S^{11},\\
 \mathsf{Discard}&=\mathsf{MP}_1;\mathsf{AX}_{1,0,0};
                         \mathsf S^{10};\mathsf{MP}_0.
\end{align*}
Their complete contracts are
\begin{align*}
 \mathsf{Push}\langle t,w,d,h\rangle&=\langle0,[t;w],0,0\rangle,\\
 \mathsf{Pop}\langle t,[a;w],d,h\rangle
 &=\begin{cases}\mathsf{halt},&t=1,\\
                 \langle t,w,0,a\rangle,&t\ne1,
   \end{cases}\\
 \mathsf{Discard}\langle t,w,0,0\rangle&=\langle0,w,0,0\rangle.
\end{align*}
Push's final MP fails because it compares residues zero and two.
Pop's final Shift tests exactly the saved top.  Discard's first MP
fails since $I(t,1)>0$; the AX identity then leaves $U(1,0)$ in
the top and zero in the cache, so its last MP fails too.  These arguments
include $t=w=0$ and the alias $[0;0]=0$.

Let $\mathsf{Offset}_0$ be empty and
$\mathsf{Offset}_u=\mathsf{MP}_u$ for $u>0$.  On a store with
$t=d=h=0$, it sets $d$ to $u$ and preserves the stack.  Then
\[
 \mathsf{Gen}_i=\mathsf{Push};\mathsf{Offset}_i;\mathsf S;
                  \mathsf{AX}_{i,0,0};\mathsf S^{10}
\]
takes $\langle a,w,0,0\rangle$ to
$\langle U(i,a),w,0,0\rangle$, including $a=1$.
The Shift preceding AX tests cache zero; AX clears the retrieved
$a$ before its subsequent Shift tests.

For $1\le j\le7$, the word
\[
 \mathsf{Discard};\mathsf{AX}_{A,D,H};\mathsf S^{j+1};
                              \mathsf{Pop}^{9-j}
\]
returns $\langle q_j,w,0,0\rangle$ from
$\langle t,w,0,0\rangle$.  After the Shifts, the top is $q_j$,
the cache is $q_{j+1}$, and the stack is
$[q_{j+2},\ldots,q_9,0;w]$.  There are $9-j$ displayed entries to
consume; each Pop preserves the safe top $q_j\ne1$, and the last
consumed entry is zero.  Thus the diagonal selector is explicitly
\[
 \mathsf{Discard};\mathsf{AX}_{B,0,0};\mathsf S^2;\mathsf{Pop}^{8}.
\]
It works at $B=0$ and with a failed gate guard.  For $u\le a$, the
propositional selector
\[
 \mathsf{Discard};\mathsf{Offset}_u;
 \mathsf{AX}_{b,a-u,c};\mathsf S^9;\mathsf{Pop}
\]
returns $\rho(a,b,c,u)$.  The gate selector
$\mathsf{Discard};\mathsf{AX}_{c,b,a};\mathsf{Pop}^{10}$ returns
$K(a,b,c)$.  All these selectors preserve the arbitrary underlying
stack and clear $d,h$.

For modus ponens the exact suffix is
\[
 \mathsf{MPSuffix}_b=\mathsf S;\mathsf S;\mathsf{MP}_b;\mathsf S.
\]
On $\langle0,[a,I(a,b),0;w],0,0\rangle$ it halts if $a=1$ or
$b=1$ and otherwise returns $\langle b,w,0,0\rangle$.
The first Shift retrieves $a$; the second tests $a$ and puts
$I(a,b)$ in the cache; MP puts $b$ in the cache; the final Shift
tests $b$ and consumes the dummy zero.  That zero is necessary to
preserve $w$.  The same calculation applies when $a=b=w=0$.

\begin{lemma}[Uniform compiled refutations]
\label{cal:word-lemma}
There is a primitive-recursive compiler $p\mapsto W_p$ such that,
for every $\mathcal C_\Gamma$ proof $p$ of $B$ and every $t,w$,
$W_p$ from $\langle t,w,0,0\rangle$ halts or returns
$\langle B,w,0,0\rangle$.  A proof of $1$ yields a command word
that halts from every incoming store.  PA proves these assertions
about the coded words and derivations.
\end{lemma}

\begin{proof}
Compile axioms by the selectors and generalization by
$W_{p_0};\mathsf{Gen}_i$.  For an MP node with minor proof
$p_0$ of $a$ and major proof $p_1$ of $I(a,b)$, use
\[
 \mathsf{Discard};\mathsf{Push};W_{p_1};\mathsf{Push};
 W_{p_0};\mathsf{Push};\mathsf{MPSuffix}_b.
\]
The first two words introduce the dummy zero, and the premise
computations and Pushes establish the suffix's three-entry
precondition.  Induction on the proof gives the claimed invariant;
an early halt remains a halt after extending the word.  If $p$
proves $1$, then
\[
            \mathsf{AX}_{0,0,0};W_p;\mathsf{Pop}
\]
halts from every incoming store.  AX first clears $d,h$; the
compiler accepts its arbitrary resulting top and stack; and the
final Pop detects top one if execution has not halted earlier.
The proof consists of the displayed total projection identities
and induction on a finite derivation, so PA verifies the compiler.
\end{proof}

\subsection{Canonical stack annotations and the exhaustive search}
\label{cal:scheduler}

For $w>0$ one has $\operatorname{tl}(w)<w$.  Thus
\[
 \delta(w)=\min\{k\le w:\operatorname{tl}^{,k}(w)=0\}
\]
is a total primitive-recursive function, with its defining property
provable in PA.  Annotate the $\delta(w)$ live heads by proofs, and
annotate both $t$ and $h$.  The tail beyond those live heads supplies
the fixed truth proof.  Pushing $a$ prepends its proof unless
$a=w=0$, in which case the canonical live stack remains empty.
Popping zero supplies truth and leaves the empty annotation.  This
is the exact treatment of the Cantor zero alias.

With either numerical proofs or their PA translations as annotations,
AX preserves the invariant: its saved top is known, its heads are
axioms, $d\le d+x$ justifies its restricted $\rho$ instance, and
$U(d,h)$ follows by generalization.  A successful MP uses the
annotated $t$ and $I(t,d+x)=h$ to prove $d+x$; a failed MP returns
truth.  Shift moves annotations or halts with a proof of $1$ in
the cache.  The arbitrary offset $d$ has no theoremhood requirement.
Every update is primitive recursive, and PA proves preservation
through every coded finite command sequence.

The retained scheduler exhaustively decodes candidate command words.
Its inherited coding and padding lemma makes every finite word
available at a candidate number above any fixed starting cutoff,
including the clean seed's counter $2$.  Its store is carried between
trials.  Lemma~\ref{cal:word-lemma}'s universal incoming-store
property is therefore essential: a compiled refutation halts when
its trial is reached, unless an earlier trial has already halted.
The candidate induction, word decoding, and terminating arithmetic
macro routines have primitive-recursive witnesses with PA-verifiable
correctness, as in the inherited PA196 source compiler.
Conversely, the annotation invariant propagates from the clean store
through all prior trials and turns any halting command into a
refutation.  Together with Lemma~\ref{cal:logical-lemma}, this proves
in PA that the Combined source halts exactly when PA has a refutation.

\subsection{From the source to the exact 194-state table}
\label{cal:physical}

The finite source check establishes equality of the entire
3,649-instruction control-flow map with the established 195C source,
under primitive-site ordinal identification.  The comparison includes
every opcode, register operand, and successor, including both decrement
branches.  It is stronger than equality of the flattened AX word.
The first 3,634 instructions match the original PA194 source except
for exit retargetings at ordinals $48,3553,3633$.  The new MP-success
tail begins at $3634$, the return to fetch is $3637$, and Shift
begins at $3638$.

The MP-success loop transfers $(D,H)=(y,0)$ to $(0,y)$ by a
decrement, increment, and backward jump, then returns to fetch.
The Shift tail clears $T$, copies $H$ to $T$ and scratch $X$,
and restores $H$ from $X$.  With the inherited $X=0$ boundary
invariant, it reaches $(T,H,X)=(h,h,0)$ in exactly
$2t+7h+3$ primitive steps; two more instructions set $X=1$ and
enter the unchanged guard.  Induction on the respective unary
counters proves these contracts for all natural store values.
No AX formula, parameter assignment, or stack order changes.

Fresh physical checks establish all 4,749 source outcomes against
the final 194-state table.  In particular, the independent literal-tape
checker rederives four generic list-sweep lemmas from that table and
uses no arithmetic-register oracle.  Its frames quantify over arbitrary
natural registers, an affine arbitrary gap, and an arbitrary finite
nonempty spare list.  Each nonhalting source step has a finite physical
run to its exact next complete frame; every duration is positive.
Increment appends a spare block, while decrement preserves the file
length, so the nonempty-spare invariant persists from initialization.
The complete blank run reaches the clean site-$25$ frame after
57,922 transitions, at native state $84$, head $-2$, logical registers
$(2,0,0,0,0,0,0,0,0)$, and 26 spare registers of value one.
The full tape is compared with the specified frame.

These facts admit a specific arithmetic internalization.  A literal
symbolic transition is a finite table lookup.  A compressed unary scan
is proved by induction on its run length.  Each generic list sweep
is proved by induction on its blocks and finite list.  The remaining
symbolic operations only reexpress the same tape.  Their trace
constructors, frame encodings, and concatenations are primitive
recursive, with correctness proved by these inductions in PA.
The machine-specific part is a finite collection of closed table
equalities and symbolic traces.  The Python checks verify that finite
data; they do not themselves constitute a formal PA proof checker.

\begin{theorem}[PA calibration of the literal PA194C table]
\label{cal:main-theorem}
Let $H_{194C}$ be the standard arithmetical sentence asserting a finite
halting trace of \file{construction/output/pa194C.tm} from its designated
blank start.  Then
\[
 \boxed{\PA\vdash
       \bigl(\neg H_{194C}\leftrightarrow\Con(\PA)\bigr).}
\]
\end{theorem}

\begin{proof}
A PA refutation gives a numerical refutation by
Lemma~\ref{cal:logical-lemma}, then a uniformly halting command word
by Lemma~\ref{cal:word-lemma}.  The exhaustive scheduler and arithmetic
macro contracts give a finite source halt.  Concatenate the exact
startup and the physical source-step traces to obtain $H_{194C}$.
All transformations and correctness inductions are available in PA.

For the reverse direction, reason in PA under the temporary assumption
$\Con(\PA)$.  The translated annotation invariant precludes a
halting command and supplies an annotated successor store at every
command boundary.  For any $T$, primitive recursion constructs
$T+1$ nonhalting command segments and their positive source and
physical simulations.  Including startup, this is a physical run
longer than $T$ with no halt before a subsequent nonhalt frame.
A purported first halting trace of length $T$ contradicts determinism:
with partial semantics it cannot be extended beyond the halt, and
with absorbing halt semantics it cannot subsequently reach a
nonhalt frame.  A halt within startup is excluded by its exact
finite trace.  No inverse decoding of an arbitrary intermediate
tape or unbounded search for a return time is needed.
Thus PA proves $\Con(\PA)\to\neg H_{194C}$.  Classical contraposition
of the already proved refutation-to-halt implication gives the other
direction.
\end{proof}

The full-induction supplier in \eqref{cal:comprehension} explains why
one cannot replace the searched calculus by the concise iPRA calculus
and retain this conclusion without a new calibration.  The iPRA
project contributes the strategy of replacing expensive primitive
schemas by smaller derived interfaces; it does not provide an
identification of $\Con(\mathrm{iPRA})$ with $\Con(\PA)$.
Likewise, the existing conditional iPRA machine endpoints in the
repository do not prove the two machine-specific reductions needed
for an unconditional iPRA theorem.  The theorem above has the
ordinary mathematical PA-provability scope stated at the start.
